\section{Глава~\ref{sec:eetypes}. Нейроны как приборы}

Режимы работы отдельных клеток измерены прямо, током через электрод и записью напряжения; математика интегратора и деления на классы --- классическая теория, а то, что мозг пользуется перестройкой режимов для вычислений, остаётся гипотезой.

{\footnotesize
\setlength{\tabcolsep}{3pt}\renewcommand{\arraystretch}{1.15}
\begin{longtable}{>{\raggedright}p{0.5cm}>{\raggedright}p{4.0cm}>{\raggedright}p{5.6cm}>{\raggedright}p{2.3cm}>{\raggedright\arraybackslash}p{1.7cm}}
\toprule
№ & Утверждение & Опыт и что измерено & Источник & Статус \\
\midrule\endfirsthead
\toprule
№ & Утверждение & Опыт и что измерено & Источник & Статус \\
\midrule\endhead
1 & Резонатор: медленный ток, противодействующий изменению напряжения, делает нейрон полосовым фильтром с максимумом на единицах--десятках герц (раздел~\ref{sec:resonator}) & В срезах в \nt{GLUT}-нейроны вводили синусоидальный ток с плавно растущей частотой и мерили импеданс $|Z(f)|$: максимум на $2$--$5$~Гц при $33^\circ$C (около $7$~Гц при $38^\circ$C); резонанс создают медленные калиевый и катионный токи и усиливает постоянный натриевый ток. Обзор даёт тот же приём для многих классов клеток & \cite{HuVervaekeStorm2002,HutcheonYarom2000} & измерено \\
2 & Медленный ток ведёт себя как индуктивность, вместе с ёмкостью мембраны --- колебательный контур & Подпороговый отклик аксона на ток измерен и описан линеаризованными уравнениями проводимостей; эквивалентная схема содержит «феноменологическую» индуктивность, величина которой зависит от напряжения & \cite{Mauro1970} & теория \\
3 & Постоянная времени группы с обратной связью $\tau_{\text{эфф}}=\tau/(1-w)$~\eqref{eq:perfectint}; при $\tau=0.1$~с и $w=0.995$ --- $20$~с & Выводится в главе из линейного уравнения группы; как механизм удержания положения глаза предложено в теоретической работе & вывод в главе; \cite{Seung1996} & теория \\
4 & Интегратор положения глаза держит вход сетью, а не свойством отдельной клетки & Внутриклеточная запись у рыбы в узле, держащем положение глаза: после быстрого поворота частота нейрона ступенькой меняется и держится; короткий ток в одну клетку вызывает лишь кратковременный сдвиг, а ступеньки напряжения сохраняются и ниже порога --- их поддерживают синаптические входы & \cite{Aksay2001} & измерено \\
5 & Интегратору без утечки нужна постоянная подстройка обучением; при расстройке он либо забывает, либо уходит в насыщение & Рыбу обучали зрительной обратной связью, пропорциональной $\pm$ положению глаза: удержание становилось неустойчивым или с утечкой, постоянная времени нейронов падала больше чем на порядок, до $<1$~с; обычная зрительная обратная связь постепенно возвращала устойчивость & \cite{Major2004} & измерено \\
6 & Бистабильный нейрон: короткий вход включает длительное плато, торможение выключает; свойство включается серотонином (раздел~\ref{sec:bistable}) & Внутриклеточная запись \nt{ACET}-нейронов, управляющих мышцами, у кошки: короткое синаптическое возбуждение переводит нейрон в длительную работу, короткое торможение сбрасывает; у спинальной кошки состояние появляется после введения предшественника серотонина & \cite{Hounsgaard1988} & измерено \\
7 & Два класса: интеграторы (частота растёт от нуля) и резонаторы (при пороге сразу конечная частота) & Классы выведены из теории бифуркаций. Опыт: в срезах коры \nt{GLUT}-нейроны начинают работать со сколь угодно малой частотой (тип~1), быстрые \nt{GABA}-нейроны --- скачком с конечной частоты (тип~2) & \cite{Izhikevich2007,Tateno2004} & измерено \\
8 & Один нейрон --- интегратор или детектор совпадений: торможение укорачивает окно сложения & В срезах гиппокампа окно, в котором возбуждающие входы складываются до порога, короче $2$~мс; его укорачивает прямое торможение, приходящее сразу за возбуждением; в дендритах, где торможение слабее, окно шире & \cite{Pouille2001} & измерено \\
9 & Линия задержки из аксонов (таблица~\ref{tab:eetypes}) & У сипухи измерены задержки в аксонах, идущих к детекторам совпадений: разная длина аксона даёт разную задержку, и детекторы настроены на разность времён прихода звука & \cite{Carr1990} & измерено \\
10 & Ритмоводитель в бодрствовании и молчаливая клетка во сне & \nt{SERO}-нейроны работают почти регулярно днём, замедляются в медленном сне и почти молчат в быстром (подробнее --- глава~\ref{sec:sleep}) & \cite{JacobsAzmitia1992,McGintyHarper1976} & измерено \\
11 & Прибор задают ионные каналы и широковещательные каналы, которые их подстраивают (утверждение~\ref{prop:eetypes}) & Показано на малых сетях: вещества-регуляторы меняют собственные свойства нейронов и силу синапсов, так что одна и та же схема выдаёт разные ритмы & \cite{Marder2012} & измерено \\
12 & Мозг использует перестройку режимов для вычислений (утверждение~\ref{prop:eetypes}) & Прямого опыта, где перестройка режима клетки была бы нужна для решения задачи, нет; есть лишь опыты, где регуляторы меняют выход схемы & \cite{Marder2012} & гипотеза \\
\bottomrule
\end{longtable}}
